% Источники: exam/materials/моделирование.pdf, раздел 1; https://github.com/HumsterProgrammer/iu9-notes/blob/main/8sem/Моделирование/Моделирование_лекция-01_09-02-26.md.
\section{Математическое моделирование. Математическая модель. Испытания, эксперимент, опыт.}

\textbf{Математическое моделирование} — метод исследования объектов и процессов реального мира с помощью их приближенных описаний с помощью формул (равенств, неравенств, уравнений, логических структур). Такое описание называют математическими моделями.

\textbf{Модель} — идеализированное представление, обладающее достаточной для решения задачи моделирования степени адекватности, исследуемой системы (объекта, процесса).

\textbf{Процесс создания математической модели:}
\begin{enumerate}
  \item построение математической модели;
  \item постановка исследования и решение вычислительных задач, реализующих модель;
  \item проверка качества моделей на практике, модификация моделей.
\end{enumerate}

Задача моделирования заключается в создании (построении через процедуру формализации) модели сложной системы с последующим построением и проведением эксперимента над моделью и анализом результатов.

\textbf{Испытание} — это воспроизведение какого-либо явления или наблюдения за динамикой исследования исследуемого объекта в естественных условиях среды, реже в искусственных. Испытание включает серию экспериментов.

\textbf{Эксперимент} — система операций, воздействий и (или) наблюдений, направленных на получение информации об объекте при исследовательских испытаниях.

\textbf{Опыт} — воспроизведение исследуемого явления в определенных условиях проведения эксперимента при возможности регистрации его результатов.

\clearpage
